\bookchapter{Kapitel 3: Das Handy an die „Lebensader“ anschließen}

\chapterintro{Null. Nicht orientierungslos im Netz — WLAN und Mobilfunk richtig unterscheiden}

\booksection{Null. Nicht orientierungslos im Netz — WLAN und Mobilfunk richtig unterscheiden}

WLAN ist verbunden, aber Seiten laden nicht\textbf{— weil der Router offline ist und das Smartphone nicht automatisch auf Mobilfunk umschaltet. Um das gründlich zu verstehen, brauchen wir ein Modell von „Straße und Fahrzeug“.}:

\begin{tcbinfo}{Kernanalogie: Straßen und Fahrzeuge}
„Straße“\textbf{ist der Kanal zwischen deinem Smartphone und dem Internet, die}„Fahrzeuge”\textbf{,}sind die Fahrzeuge, die darauf unterwegs sind. WLAN und Mobilfunk sind zwei verschiedene Straßen, doch die Daten (Fahrzeuge) sind dieselben.

\end{tcbinfo}

\booksubsection{,}

WLAN ist im Wesentlichen eine\textbf{— wie eine drahtlos verlegte Straße vor deinem Haus, die die Daten vom Router zum Smartphone bringt. Wenn aber die Breitbandrechnung nicht bezahlt ist oder das Glasfaserkabel gerissen ist, ist die Straße zwar noch da, aber die}„Autobahn“

\booksubsection{Kosten und Szenarien: feste Monatsrate vs. mengenabhängige Abrechnung}

WLAN ist meist eine einmalige Investition (Breitband-Monatsgebühr) — solange der Strom nicht ausfällt, kann die ganze Familie unbegrenzt nutzen. Ideal für datenintensive Szenarien wie Filme herunterladen oder Spielen. Mobilfunk wird nach Verbrauch abgerechnet, bei Überschreitung des Tarifs kann gedrosselt oder extra abgerechnet werden — besser für leichte Aufgaben in der U-Bahn oder auf Geschäftsreise.

\begin{tcbpractice}{Probier es aus: Prüfe deinen Mobilfunktarif}
Öffne auf dem Smartphone [Einstellungen] [Mobilfunknetz] [Datenverwaltung] und schau, wie viel Datenvolumen du diesen Monat noch hast und welche Apps am meisten verbrauchen. Wer Bescheid weiß, hat unterwegs keinen Stress.

\end{tcbpractice}

\booksubsection{Sicherheitsaspekt: öffentliches WLAN ist ein verstecktes „Angelgewässer“}

Hinsichtlich der Sicherheit sind Mobilfunkdaten in der Regel sicherer als öffentliches WLAN. Mobilfunkdaten werden vom Betreiber verschlüsselt, während öffentliches WLAN tatsächlich von Hackern als „gefälschter Hotspot“ aufgebaut werden kann, um unverschlüsselten Datenverkehr abzufangen. Hier aber genau unterscheiden: In einer Zeit, in der HTTPS allgegenwärtig ist, kann ein Angreifer nur\textbf{erfahren, welche Websites du besuchst}(Domain), sowie\textbf{unverschlüsselte HTTP-Inhalte}; dagegen können über HTTPS verschlüsselte Daten wie WeChat-Nachrichten oder die Anmelde-Passwörter der Bank-App von Angreifern\textbf{nicht direkt mitlesbar}– es sei denn, es kommt ein aufwendiger SSL-Stripping-Angriff zum Einsatz und du ignorierst die Zertifikatswarnung. Trotzdem gilt: Für Zahlungen unterwegs ist der eigene Mobilfunkkanal die sicherste Wahl.

\textbf{Sicherheitshinweis}\par

Verbinde dich an Flughäfen, in Einkaufszentren oder Cafés nicht mit kostenlosen WLANs mit verdächtigen Namen (wie „Free\_WiFi\_2024“). Und selbst wenn du verbunden bist, melde dich dort niemals in einer Bank-App an oder gib Passwörter ein.

\booksubsection{Am Handy ablesen: die „Geheimsprache“ der Statusleiste}

3. Sicherheitsaspekte: Öffentliches WLAN als heimische „Angelseite“\textbf{Pfeil-Symbol (Doppelpfeil)}– dann surfst du gerade über diesen Weg; leuchten beide Symbole, greift die Netzwerk-Beschleunigung und das Handy nutzt beide Kanäle gleichzeitig.

iPhone-Nutzer haben es einfacher: Leuchtet das WLAN-Symbol, ist WLAN aktiv; leuchtet es nicht, gilt das Balkensymbol oben rechts für Mobilfunk. Manche Android-Handys zeigen in der Statusleiste zusätzlich „5G“ oder „4G“ und verraten dir damit direkt den aktuellen Mobilfunkstandard.

\booksection{Mobilfunknetz (4G/5G) und Datenvolumen}

Das Mobilfunknetz ist die Basis, um überall und jederzeit online zu gehen. Von 2G bis 5G wuchs das Tempo um das Tausendfache – das Prinzip blieb aber im Kern dasselbe.

\booksubsection{Netzstandards:}

\begin{itemize}
  \item \textbf{4G LTE}: Derzeit die größte Abdeckung, ausreichende Geschwindigkeit (theoretisch über 100Mbps+), vergleichsweise geringer Stromverbrauch.
  \item \textbf{5G}: Schnell (theoretisch bis 10Gbps), aber die Abdeckung wird noch ausgebaut, höherer Stromverbrauch. Geeignet für große Downloads und hochauflösende Videos.
  \item \textbf{: SA ist Stand-Alone, NSA basiert auf 4G-Basisstationen. SA bietet das bessere Erlebnis, erfordert aber die Unterstützung von Betreiber und Smartphone.}: SA ist eigenständiges Netzwerken (Standalone), NSA stützt sich auf 4G-Sendemasten. SA bietet das bessere Erlebnis, setzt aber voraus, dass Netzbetreiber und Handy es unterstützen.
\end{itemize}

\booksubsection{Datenmanagement:}

\begin{itemize}
  \item \textbf{Was ist Datenvolumen?}: Die Datenmenge, die beim mobilen Surfen anfällt. Videos und Bildfeeds verbrauchen am meisten, Text und Sprachnachrichten sind sparsam.
  \item \textbf{Spartipp}: In den Einstellungen den „Datensparmodus“ aktivieren und App-Updates bzw. HD-Videos aufs WLAN verschieben.
  \item \textbf{Balken $\neq$ Tempo}: Starker Empfang bedeutet nur eine gute Verbindung – das Tempo hängt davon ab, wie ausgelastet die Funkzelle ist und wie weit du von ihr entfernt bist.
\end{itemize}

\textbf{Wichtige Warnung: Vor der Reise „Datenroaming“ ausschalten!}\par

Mobile Daten im Ausland können himmlische Rechnungen verursachen! Schalte vor der Abreise in den Einstellungen das Datenroaming aus oder kaufe vorab eine lokale SIM/eSIM oder einen internationalen Roaming-Tarif.

\booksection{WLAN: Verbinden, Teilen, Sicherheit}

Wi-Fi 4

\booksubsection{Mit WLAN verbinden:}

: Aktiviere in den Einstellungen den „Daten sparren“-Modus und aktualisiere Apps oder schaue hochauflösende Videos nur unter WLAN.

\booksubsection{Passwort teilen:}

Android und iOS können das WLAN-Passwort schnell an Freunde weitergeben. Android erzeugt einen QR-Code, iOS teilt automatisch über Nahbereich.

\booksubsection{Wi-Fi 5}

WLAN-Standards werden von der IEEE festgelegt, jede Generation bringt deutliche Verbesserungen:

\begin{center}
\begin{tabularx}{\linewidth}{|X|X|X|X|X|}
\hline
\textbf{Standard} & \textbf{Protokoll} & \textbf{Frequenzband} & \textbf{Höchstrate} & \textbf{Wichtige Warnung: Schalte vor der Auslandsreise Datenroaming aus!} \tabularnewline
\hline
Wi-Fi 4 & 802.11n & 2.4/5GHz & 600Mbps & Alter Standard, noch ordentliche Abdeckung \tabularnewline
Wi-Fi 6E & 802.11ac & 5GHz & 3.5Gbps & Weitestgehend verbreitet, Tempo ausreichend \tabularnewline
Wi-Fi 7 & 802.11ax & 2.4/5GHz & 9.6Gbps & Stark bei vielen Geräten, niedrige Latenz \tabularnewline
: Führt MLO (Multi-Link Operation) ein, kann mehrere Frequenzbänder gleichzeitig nutzen und die Latenz deutlich senken. 320MHz ultrabreite Kanäle sorgen für einen Geschwindigkeitssprung. Ab 2024 unterstützen Flagship-Chips wie Snapdragon 8 Gen3/8 Elite, Dimensity 9300/9400 u. a. Wi-Fi 7. Modelle wie das Xiaomi 15, das Huawei Mate 70 und das Samsung S25 können Wi-Fi-7-Netze nutzen (in Kombination mit einem Wi-Fi-7-Router). & 802.11ax & 2.4/5/6GHz & 46Gbps & 2.4/5/6GHz \tabularnewline
: 2.4GHz durchdringt Wände besser, ist aber langsamer und störanfälliger; 5GHz ist schneller, aber durchdringt Wände schlechter. Das 6GHz-Band von Wi-Fi 6E/7 ist am schnellsten, durchdringt aber am schlechtesten — geeignet für die Nutzung im selben Raum. & 802.11be & : 2.4GHz durchdringt Wände besser, ist aber langsamer und störanfälliger; 5GHz ist schneller, aber durchdringt Wände schlechter. Das 6GHz-Band von Wi-Fi 6E/7 ist am schnellsten, durchdringt aber am schlechtesten — geeignet für die Nutzung im selben Raum. & 46 Gbit/s & Höchstgeschwindigkeit, Multi-Link-Bündelung \tabularnewline
\hline
\end{tabularx}
\end{center}

\begin{itemize}
  \item \textbf{: Erweitert Wi-Fi 6 um das 6GHz-Band, bietet mehr Kanäle und reduziert Störungen. Setzt Unterstützung auf Smartphone und Router voraus.}: OFDMA (Orthogonal Frequency Division Multiple Access) lässt viele Geräte gleichzeitig senden, MU-MIMO (Multi-User-MIMO) erhöht die Effizienz bei vielen Teilnehmern, und TWT (Target Wake Time) weckt Endgeräte gezielt auf – zusammen senken sie den Verbrauch.
  \item \textbf{: Führt MLO (Multi-Link Operation) ein, kann mehrere Frequenzbänder gleichzeitig nutzen und die Latenz deutlich senken. 320MHz ultrabreite Kanäle sorgen für einen Geschwindigkeitssprung. Ab 2024 unterstützen Flagship-Chips wie Snapdragon 8 Gen3/8 Elite, Dimensity 9300/9400 u. a. Wi-Fi 7. Modelle wie das Xiaomi 15, das Huawei Mate 70 und das Samsung S25 können Wi-Fi-7-Netze nutzen (in Kombination mit einem Wi-Fi-7-Router).}: Erweitert Wi-Fi 6 um das 6-GHz-Band – mehr Kanäle, weniger Störungen. Funktioniert nur, wenn sowohl Smartphone als auch Router es unterstützen.
  \item \textbf{Wi-Fi 7}: Führt MLO ein (Multi-Link Operation): Mehrere Bänder übertragen gleichzeitig, was die Latenz deutlich senkt, und 320-MHz-Superkanäle lassen das Tempo sprunghaft steigen. Top-Chips wie Snapdragon 8 Elite Gen 5 und Dimensity 9400 unterstützen Wi-Fi 7 bereits; Modelle wie Xiaomi 17, Huawei Mate 80 und Samsung S26 können Wi-Fi-7-Netze nutzen (mit Wi-Fi-7-Router).
  \item \textbf{2,4 GHz vs. 5 GHz}: Fuehrt MLO (Multi-Link Operation) ein, kann mehrere Frequenzbaender gleichzeitig fuer die Datenuebertragung nutzen und die Latenz drastisch senken. 320-MHz-Ultrabreitkanal fuer massive Geschwindigkeitsspruenge. Flaggschiff-Chips wie Snapdragon 8 Elite Gen5 und Dimensity 9400 unterstuetzen bereits Wi-Fi 7; Geraete wie Xiaomi 17, Huawei Mate 80 und Samsung S26 koennen Wi-Fi-7-Netze nutzen (erfordert einen Wi-Fi-7-Router).
\end{itemize}

\booksubsection{Risiken öffentlicher WLANs:}

In öffentlichen WLANs können Hacker Daten abhören oder Phishing-Hotspots aufspannen. Verhaltensregeln:

\begin{itemize}
  \item Keine Banking-Logins, keine Passworteingaben.
  \item Besser: Verbindung per VPN verschlüsseln.
  \item Fremden WLANs nicht automatisch beitreten.
  \item HTTPS-Webseiten nutzen (Schloss-Symbol in der Adressleiste).
\end{itemize}

\booksection{Bluetooth: der Allrounder unter den Funkverbindungen}

Bluetooth ist wahrscheinlich eine der am häufigsten genutzten Funktechniken deines Handys – kabellose Kopfhörer, Smartwatches, das Auto-System, Tastatur und Maus verbinden sich alle darüber.

\booksubsection{Wie die Technik funktioniert:}

Bluetooth arbeitet im 2,4-GHz-ISM-Band und nutzt Frequency-Hopping-Spread-Spectrum-Technik (FHSS), um Daten über 79 Kanäle schnell wechselnd zu verteilen und so Störungen effektiv zu vermeiden. Die Kommunikationsreichweite von Bluetooth ist in drei Klassen eingeteilt: Klasse 1 ca. 100 Meter, Klasse 2 ca. 10 Meter (bei Smartphones üblich), Klasse 3 ca. 1 Meter.

\booksubsection{Versionsentwicklung:}

\begin{itemize}
  \item \textbf{Bluetooth 4.0}: Einführung von Bluetooth Low Energy (BLE), sodass Geräte mit kleiner Batterie wie Fitnessarmbänder und Herzfrequenzgurte ebenfalls drahtlos verbunden werden können – die Akkulaufzeit stieg von Tagen auf Monate.
  \item \textbf{Bluetooth 4.2}: Verbesserte Übertragungsgeschwindigkeit und Datenschutz, Unterstützung für direkte IPv6-Verbindung.
  \item \textbf{Bluetooth 5.0}: Verdoppelte Übertragungsgeschwindigkeit (2 Mbit/s), vierfache Reichweite (theoretisch 300 Meter), achtfache Broadcast-Kapazität. Unterstützung für Dual Audio (gleichzeitiges Verbinden von zwei Kopfhörern).
  \item \textbf{Bluetooth 5.1}: Einführung der Richtungsfindung (AoA/AoD), Positionierungsgenauigkeit bis auf Zentimeterebene.
  \item \textbf{Bluetooth 5.2}: 2024–2025 ist LE Audio in die schnelle Verbreitungsphase eingetreten. Neue True-Wireless-Kopfhörer aller großen Marken (z. B. Huawei FreeBuds Pro 4, Xiaomi Buds 5 Pro, Samsung Galaxy Buds3) unterstützen LE Audio und LC3-Codierung und bieten geringere Latenz sowie bessere Tonqualität. Beim Kauf neuer Kopfhörer wird empfohlen, Modelle mit LE-Audio-Unterstützung vorzuziehen.
  \item \textbf{Bluetooth 5.3/5.4}: Weitere Optimierung von Verbrauch und Verbindungsstabilität, verbessertes periodisches Advertising.
\end{itemize}

\textbf{Audioübertragung}: Kabellose Kopfhörer, Auto-Soundsysteme. LE Audio bei Bluetooth 5.2 bringt sowohl Tonqualität als auch Latenz auf ein neues Niveau.

\booksubsection{Anwendungsfälle:}

\begin{itemize}
  \item \textbf{Audioübertragung}: 2025-2026 befindet sich LE Audio in der raschen Massenverbreitung. Neue True-Wireless-Kopfhorrer aller Marken (z. B. Huawei FreeBuds Pro 5, Xiaomi Buds 6 Pro, Samsung Galaxy Buds3 Pro) unterstuetzen bereits LE Audio und LC3-Codierung mit geringerer Latenz und besserer Klangqualitaet. Beim Kauf neuer Kopfhoerer wird empfohlen, LE-Audio-Unterstuetzung vorzuziehen.
  \item \textbf{Datenübertragung}: Dateien und Fotos über kurze Distanz senden. Zwar langsamer als Wi-Fi Direct, aber dafür einfach.
  \item \textbf{Gerätepaarung}: Smartwatches, Fitnessarmbänder, Tastaturen, Mäuse, Gamecontroller.
  \item \textbf{Standortdienste}: Bluetooth-Beacons für Indoor-Navigation und Gegenstandsverfolgung (z. B. Apple AirTag).
  \item \textbf{Smart Home}: Bluetooth-Mesh-Netzwerke zur Steuerung von Glühbirnen, Türschlössern usw.
\end{itemize}

\booksection{Vier. Infrarot-Fernbedienung: Das Smartphone als Universalfernbedienung}

Die Infrarot-Fernbedienung ist eine der ältesten drahtlosen Steuerungstechnologien, bleibt aber am Smartphone weiterhin praktisch. Einige Android-Smartphones (z. B. Xiaomi und bestimmte Huawei-Modelle) verfügen über einen integrierten Infrarots-Sender und können Fernseher, Klimaanlagen, Set-Top-Boxen und andere Haushaltsgeräte fernsteuern.

\booksubsection{Funktionsprinzip:}

Die Infrarot-Fernbedienung sendet über eine Infrarot-LED Lichtpulse einer bestimmten Frequenz (meist 38 kHz); der Infrarot-Empfänger des Haushaltsgeräts dekodiert diese Pulse und führt die entsprechende Aktion aus.

\booksubsection{Einschränkungen:}

\begin{itemize}
  \item \textbf{Richtungsabhängigkeit}: Infrarot überträgt geradlinig und muss direkt auf das Gerät gerichtet werden; keine Übertragung durch Wände.
  \item \textbf{Kurze Reichweite}Suche mit einem anderen Gerät (z. B. Computer) nach diesem WLAN und verbinde dich.
  \item \textbf{Keine Apple-Unterstützung}: Das iPhone hat keinen Infrarot-Sender und kann diese Funktion nicht nutzen.
\end{itemize}

\booksection{Fünf. NFC: Die Magie des Berührens}

NFC basiert auf der RFID-Technik und unterstützt drei Betriebsmodi:

\booksubsection{Kommunikationsprinzip und Modi:}

NFC ist aus der RFID-Technik hervorgegangen und unterstützt drei Betriebsmodi:

\begin{itemize}
  \item \textbf{Lese-/Schreibmodus}: Das Smartphone liest die Informationen in einem NFC-Tag, etwa eine NFC-Karte aktivieren, die eine Aktion auslöst.
  \item \textbf{Punkt-zu-Punkt-Modus}: Zwei NFC-Geräte übertragen Daten gegenseitig, z. B. Android Beam (eingestellt).
  \item \textbf{Kartenemulationsmodus}: Das Smartphone emuliert eine Smartcard – das ist der am häufigsten verwendete Modus. Transit-Karten, Zugangskarten und Bankkarten nutzen alle diesen Modus.
\end{itemize}

\booksubsection{Sicherheitsmechanismen:}

\begin{itemize}
  \item \textbf{Extrem kurze Reichweite}: Die Kommunikationsreichweite von unter 10 cm ist selbst eine physische Sicherheitsgarantie; ein Abhören aus der Ferne ist nahezu unmöglich.
  \item \textbf{Sicherheitschip}: Das Smartphone verfügt über einen unabhängigen Sicherheitschip (SE) oder nutzt HCE (Host Card Emulation). Die Zahlungsschlüssel werden im sicheren Bereich gespeichert und können selbst bei gerootetem Smartphone nicht extrahiert werden.
  \item \textbf{Fingerabdruck-/Gesichtsverifizierung}: Bei Zahlungen ist eine biometrische Bestätigung erforderlich, sodass selbst bei Verlust des Smartphones kein Missbrauch möglich ist.
  \item \textbf{Tokenisierung}: Apple Pay, Huawei Pay u. a. übertragen nicht die echte Kartennummer, sondern verwenden Einmal-Token. Selbst bei Abfangen können diese nicht wiederverwendet werden.
\end{itemize}

\booksubsection{Typische Anwendungen:}

\begin{itemize}
  \item \textbf{Transit-/Zugangskarte}: Physische Karteninformationen aufs Smartphone schreiben und einfach das Smartphone ans Lesegerät halten.
  \item \textbf{Mobile Zahlung}: Apple Pay, Huawei Pay usw. – die Bankkarte ans Smartphone binden und per Berührung bezahlen.
  \item \textbf{Gerätepaarung}: NFC-Berührung zum schnellen Paaren von Bluetooth-Kopfhörern und Lautsprechern.
  \item \textbf{Identitätsverifizierung}: NFC-Ausweislesung, elektronischer Führerschein-Nachweis.
\end{itemize}

\booksection{Sechs. Persönlicher Hotspot: Das Smartphone als WLAN-Router}

Wenn du selbst viel Datenvolumen hast, kannst du den Hotspot für Tablet oder Computer aktivieren.

\begin{itemize}
  \item \textbf{Aktivierung}: Einstellungen > Persönlicher Hotspot/Netzwerkfreigabe, einschalten.
  \item \textbf{: Apple Pay, Huawei Pay \& Co. übertragen keine echte Kartennummer, sondern Einmal-Tokens – selbst abgefangen lassen sie sich nicht wiederverwenden.}: Sehr strom- und datenintensiv – ein komplexes Passwort festlegen und die Anzahl verbundener Geräte begrenzen.
  \item \textbf{USB-Freigabe}: Mit Datenkabel an den Computer anschließen – stabiler und schneller als ein drahtloser Hotspot.
\end{itemize}

\begin{tcbpractice}{Selbst ausprobieren: Hotspot einrichten und verwalten}
\begin{enumerate}
  \item Öffne [Einstellungen], finde „Persönlicher Hotspot" oder „Netzwerkfreigabe", aktiviere ihn und lege Namen und Passwort fest.
  \item Suche mit einem anderen Gerät (z. B. Computer) nach diesem WLAN und verbinde dich.
  \item Prüfe auf dem Smartphone, ob du verbundene Geräte verwalten oder unbekannte Geräte entfernen kannst.
\end{enumerate}

\end{tcbpractice}
